% Generated draft: requires mathematical and source review.
\documentclass[11pt]{article}
\usepackage[margin=25mm]{geometry}
\usepackage{fontspec}
\setmainfont{Latin Modern Roman}
\usepackage{amsmath,amssymb,mathtools,unicode-math}
\setmathfont{Latin Modern Math}
\usepackage{enumitem,needspace,xurl,hyperref}
\setlength{\parindent}{0pt}
\setlength{\parskip}{5pt}
\setlength{\emergencystretch}{3em}
\newcommand{\N}{\mathbb{N}}
\newcommand{\Z}{\mathbb{Z}}
\newcommand{\Q}{\mathbb{Q}}
\newcommand{\R}{\mathbb{R}}
\newcommand{\C}{\mathbb{C}}
\newcommand{\F}{\mathbb{F}}
\newcommand{\A}{\mathbb{A}}
\newcommand{\PP}{\mathbb{P}}
\newcommand{\E}{\mathbb{E}}
\newcommand{\eps}{\varepsilon}
\newcommand{\dd}{\,\mathrm{d}}
\begin{document}
\section*{UnsolvedMath ID: 30005153}
\textbf{Source problem code:} catalogue entry
\subsection{Definitions and mathematical statement}
The supplied target is to establish full asymptotic stability of the $\phi^4$ kink under odd perturbations, including control of the quadratic source term generated by the internal mode. The record does not specify the evolution equation, its normalization, the perturbation space, the convergence topology, or the quadratic source term. Consequently, it does not determine a unique quantified mathematical proposition.

For clarity, in a one-dimensional formulation with spatial variable $x\in\R$, an odd perturbation $p$ satisfies
\[
 p(-x)=-p(x)\qquad\text{for every }x\in\R.
\]
For a second-order evolution with perturbation data $(p_0,p_1)$, requiring both components to be odd would mean
\[
 p_j(-x)=-p_j(x)\qquad(x\in\R,\ j\in\{0,1\}).
\]
These formulas explain a possible meaning of odd perturbations; the spatial domain and the intended condition on the complete initial data are not specified in the supplied record.

In the usual terminology, a stationary kink connects distinct equilibrium states at the two spatial ends, and an internal mode is a spatially localized eigenmode of the linearization about the kink. The record supplies neither the kink profile nor the linearized operator. Its reference to a quadratic source term indicates forcing of second order in the internal-mode amplitude, but gives no formula or requested estimate.

\textbf{Unresolved formulation.} To turn the supplied target into a precise stability statement, one must specify the equation and spatial domain; the kink and admissible initial data; the smallness and regularity assumptions; the meaning of full asymptotic stability, including convergence topology and any allowed translation or velocity parameters; and the quadratic forcing term and the estimate required to control it. None of these choices is fixed here.

There is also a substantive conflict within the supplied evidence. The extracted statement and original passage identify full asymptotic stability for odd perturbations as the open target. The supplied dated literature assessment instead reports that asymptotic stability for odd perturbations in the energy space has been proved, while the full problem for general perturbations remains unresolved. The record does not explain whether its word ``full'' imposes a stronger conclusion in the odd class. Thus neither replacing odd perturbations by general perturbations nor selecting a stronger odd-perturbation theorem is justified by the supplied material.
\subsection{Short English statement}
The record requests asymptotic stability of the $\phi^4$ kink for odd perturbations, with control of forcing from the internal mode. Essential analytic specifications are missing, and the supplied literature assessment reports an odd-perturbation result already proved. The intended remaining theorem therefore requires clarification and human mathematical review.
\subsection{Sources}
\begin{itemize}
\item Supplied local problem registry: permanent UnsolvedMath ID \texttt{30005153}; exact source problem code \texttt{catalogue entry}; title: Asymptotic Stability of the $\phi^4$ Kink. No individual proposer is supplied. The recorded proposed year is 2022.
\item Supplied report citation: \emph{Nonlinear Waves and Dispersive Equations (2023)}, Oberwolfach Reports, pp.~1661--1730, DOI \texttt{10.4171/owr/2022/30}. \url{https://doi.org/10.4171/owr/2022/30}. The supplied background describes Oberwolfach Reports as documenting mathematical workshops and their open-problem discussions. No precise page or named contributor for this extracted passage is supplied.
\item Supplied literature source, labeled ``Kowalczyk--Martel--Munoz, odd perturbations'': \url{https://arxiv.org/abs/1506.07420}.
\item Supplied literature source, labeled ``Kowalczyk--Martel, odd kink dynamics'': \url{https://arxiv.org/abs/2203.04143}.
\item Supplied literature source, labeled ``OWR report'': \url{https://ems.press/content/serial-article-files/46968}.
\item Literature assessment dated exactly \texttt{2026-08-22}: the supplied curation labels the problem partially solved and reports asymptotic stability of the classical $\phi^4$ kink for odd perturbations in the energy space as proved, with full asymptotic stability for general perturbations unresolved. This is supplied dated curation, not an independently verified literature assessment. The background explicitly calls for independent verification of time-sensitive status and cited claims.
\item The registry supplies the statement-status label \texttt{corrected\_verified} and says the statement was reconstructed and checked against the official report. These are recorded provenance claims, not verification performed for this draft.
\end{itemize}
\end{document}
